% Page 27 — Nouvelle hypothèse sur le nom d'Alban

\begin{center}
  {\Large\itshape Nouvelle hypothèse sur le nom d’Alban}
\end{center}
\vspace{1.2em}

\noindent
\lettrine{A}{ucune} des hypothèses émises jusqu’ici sur l’origine du mot d’ALBAN n’est absolument satisfaisante pour l’esprit. En voici une autre.

Il y a 17 « Saint Alban » en France, mais un seul « ALBAN ».

\medskip

\noindent
\lettrine{P}{ar} ailleurs,

\medskip

\hspace*{1.5em}D’après le savant Ducange \textit{albani} (pluriel) signifie : colons.

\hspace*{1.5em}D’après une brochure de l’Ambassade d’Albanie à Paris, c’est une tribu illyrienne, « les Albans » qui auraient donné son nom à la ville d’ALBANOPOLIS (ville d’Alban), plusieurs siècles avant J.C. et finalement au pays tout entier : l’ALBANIE.

\hspace*{1.5em}D’après la vie de Saint Barthélemy, le corps de cet apôtre aurait été porté à ALBANOPOLIS en Arménie, pays évangélisé par lui.

\hspace*{1.5em}À 30 kilomètres de Rome se trouve ALBANO et le lac du même nom. C’est par là que le pape a sa résidence d’été « Castelgandolfo ».

\hspace*{1.5em}D’autre part, on sait que les Gaulois, quatre siècles avant J.C. essaiment jusqu’en Arménie et en Asie Mineure. Saint Paul a une lettre aux Galates (Gaulois) Pourquoi alors ne pas penser que ces Gaulois aventuriers sont originaires de notre coin de montagne, portés par le flot gaulois de par le monde, ils y ont marqué la trace de leur passage en donnant leur nom à quelque cité ou pays. J’ai cité Albano en Italie, Albanopolis en Arménie et dans l’actuelle Albanie ?

\medskip

\noindent
\lettrine{L}{a} charte de 1231 – accord entre l’évêque d’Albi : Durand de Beaucaire et l’Abbé Adhémar du monastère Saint Benoît de Castres, portant la mention de Sancta Maria de ALBANIO pourrait donc signifier Notre Dame au pays des ALBAN, comme Notre Dame de Paris a signifié Notre Dame au pays des Parisis.

Notre Dame d’ALBAN signifierait Notre Dame des Albanais, comme Notre Dame de Paris a signifié Notre Dame des Parisiens.
